% GENERATED FILE. Edit canonical lessons, then run npm run generate:books.
% canonical-source-hash: fnv1a64:0c4ad3c64bdcf2a5
% canonical-lessons: GU-C297-cetavvu, GU-C297-savdhan, GU-C297-udhu, GU-C297-tolu, GU-C297-dhakko

\chapter{Warnings Nobody Heeded}
\label{ch:gu-warnings-nobody-heeded}
\hlchaptermodality{297}

\begin{chapteropening}
\textbf{By the end of this chapter:} \emph{I can say who warned whom and what happened.}

Complete the last lesson of chapter 297: I can say who warned whom and what happened.
\end{chapteropening}

\section[\texorpdfstring{\gu{ચેતવવું} (cetavvũ)}{ચેતવવું (cetavvũ)}]{\gu{ચેતવવું} (cetavvũ) — to warn}

\label{lesson:GU-C297-cetavvu}



\begin{quote}
Before the new one: say the Gujarati for a stitch, then the Gujarati for a condition.
\end{quote}

\subsection*{You'll want to know: \gu{ચેતવવું}}
\textbf{\gu{ચેતવવું}} — \emph{cetavvũ} — \textquotedblleft{}to warn\textquotedblright{}.

\begin{quote}
\gu{ચોકીદારે} \gu{મુસાફરોને} \gu{ચેતવ્યા} \gu{કે} \gu{પુલ} \gu{નબળો} \gu{છે}.
\end{quote}

\begin{quote}
\emph{cokīdāre musāphrone cetavyā ke pul nabḷo che.}
\end{quote}

\begin{quote}
The watchman warned the passengers that the bridge was weak.
\end{quote}

Past, with \gu{તેણે} (he or she …): \gu{તેણે} \gu{બધાને} \gu{ચેતવ્યા}, \emph{teṇe badhāne cetavyā}.

\begin{cousinweb}
\textbf{\gu{ચેતવવું}} (\emph{cetavvũ}) goes back to Sanskrit \emph{cētáyati}, \textquotedblleft{}makes attentive, observes\textquotedblright{} (CDIAL 4906): to make someone notice a danger. The person warned takes \textbf{\gu{ને}}: \textbf{\gu{મુસાફરોને} \gu{ચેતવ્યા}}.
\end{cousinweb}

\begin{grammarlens}[title={What you've built}]
One more word for telling what happened.
\end{grammarlens}

\subsection*{Your turn}
\begin{itemize}
\raggedright
  \item \emph{Say it:} \emph{cetavvũ}
  \item \emph{Say it:} \emph{cetavvũ}, once more
  \item \emph{Say it:} \emph{cokīdāre musāphrone cetavyā ke pul nabḷo che.} — The watchman warned the passengers that the bridge was weak.
  \item \emph{Recall:} say the Gujarati for a stitch, then the Gujarati for a condition, then say \emph{cetavvũ} again
\end{itemize}

\begin{tcolorbox}[breakable,colback=teal!4,colframe=teal!35!black,title={Before you move on}]
What does \gu{ચેતવવું} mean? (\textquotedblleft{}To warn\textquotedblright{}.) Say it once more.
\end{tcolorbox}
\section[\texorpdfstring{\gu{સાવધાન} (sāvdhān)}{સાવધાન (sāvdhān)}]{\gu{સાવધાન} (sāvdhān) — careful; on guard}

\label{lesson:GU-C297-savdhan}



\begin{quote}
Before the new one: say the Gujarati for a condition, then the Gujarati for to warn.
\end{quote}

\subsection*{You'll want to know: \gu{સાવધાન}}
\textbf{\gu{સાવધાન}} — \emph{sāvdhān} — \textquotedblleft{}careful; on guard\textquotedblright{}.

\begin{quote}
\gu{રાત્રે} \gu{આ} \gu{રસ્તા} \gu{પર} \gu{સાવધાન} \gu{રહો}.
\end{quote}

\begin{quote}
\emph{rātre ā rastā par sāvdhān raho.}
\end{quote}

\begin{quote}
Be careful on this road at night.
\end{quote}

\begin{cousinweb}
\textbf{\gu{સાવધાન}} (\emph{sāvdhān}) is a tatsama, Sanskrit \emph{sāvadhāna}, \textquotedblleft{}attentive, heedful\textquotedblright{} (Monier-Williams). On a sign it is a warning by itself: \textbf{\gu{સાવધાન}!}, look out!
\end{cousinweb}

\begin{grammarlens}[title={What you've built}]
One more word for telling what happened.
\end{grammarlens}

\subsection*{Your turn}
\begin{itemize}
\raggedright
  \item \emph{Say it:} \emph{sāvdhān}
  \item \emph{Say it:} \emph{sāvdhān}, once more
  \item \emph{Say it:} \emph{rātre ā rastā par sāvdhān raho.} — Be careful on this road at night.
  \item \emph{Recall:} say the Gujarati for a condition, then the Gujarati for to warn, then say \emph{sāvdhān} again
\end{itemize}

\begin{tcolorbox}[breakable,colback=teal!4,colframe=teal!35!black,title={Before you move on}]
What does \gu{સાવધાન} mean? (\textquotedblleft{}Careful; on guard\textquotedblright{}.) Say it once more.
\end{tcolorbox}
\section[\texorpdfstring{\gu{ઊંધું} (ū̃dhũ)}{ઊંધું (ū̃dhũ)}]{\gu{ઊંધું} (ū̃dhũ) — upside down}

\label{lesson:GU-C297-udhu}



\begin{quote}
Before the new one: say the Gujarati for to warn, then the Gujarati for careful.
\end{quote}

\subsection*{You'll want to know: \gu{ઊંધું}}
\textbf{\gu{ઊંધું}} — \emph{ū̃dhũ} — \textquotedblleft{}upside down\textquotedblright{}.

\begin{quote}
\gu{અકસ્માત} \gu{પછી} \gu{ગાડી} \gu{ઊંધી} \gu{પડી} \gu{હતી}.
\end{quote}

\begin{quote}
\emph{akasmāt pachī gāḍī ū̃dhī paḍī hatī.}
\end{quote}

\begin{quote}
After the accident the car lay upside down.
\end{quote}

\begin{cousinweb}
\textbf{\gu{ઊંધું}} (\emph{ū̃dhũ}) goes back to Sanskrit \emph{avamūrdha}, \textquotedblleft{}head down\textquotedblright{}, and CDIAL lists the Gujarati word, \textquotedblleft{}upside down\textquotedblright{} (CDIAL 804). It agrees like an adjective: \textbf{\gu{ઊંધી} \gu{ગાડી}}, \textbf{\gu{ઊંધું} \gu{વાસણ}}.
\end{cousinweb}

\begin{grammarlens}[title={What you've built}]
One more word for telling what happened.
\end{grammarlens}

\subsection*{Your turn}
\begin{itemize}
\raggedright
  \item \emph{Say it:} \emph{ū̃dhũ}
  \item \emph{Say it:} \emph{ū̃dhũ}, once more
  \item \emph{Say it:} \emph{akasmāt pachī gāḍī ū̃dhī paḍī hatī.} — After the accident the car lay upside down.
  \item \emph{Recall:} say the Gujarati for to warn, then the Gujarati for careful, then say \emph{ū̃dhũ} again
\end{itemize}

\begin{tcolorbox}[breakable,colback=teal!4,colframe=teal!35!black,title={Before you move on}]
What does \gu{ઊંધું} mean? (\textquotedblleft{}Upside down\textquotedblright{}.) Say it once more.
\end{tcolorbox}
\section[\texorpdfstring{\gu{ટોળું} (ṭoḷũ)}{ટોળું (ṭoḷũ)}]{\gu{ટોળું} (ṭoḷũ) — a gathered crowd; a mob}

\label{lesson:GU-C297-tolu}



\begin{quote}
Before the new one: say the Gujarati for careful, then the Gujarati for upside down.
\end{quote}

\subsection*{You'll want to know: \gu{ટોળું}}
\textbf{\gu{ટોળું}} — \emph{ṭoḷũ} — \textquotedblleft{}a gathered crowd; a mob\textquotedblright{}.

\begin{quote}
\gu{અકસ્માત} \gu{પછી} \gu{ત્યાં} \gu{લોકોનું} \gu{મોટું} \gu{ટોળું} \gu{ભેગું} \gu{થયું}.
\end{quote}

\begin{quote}
\emph{akasmāt pachī tyā̃ lokonũ moṭũ ṭoḷũ bhegũ thayũ.}
\end{quote}

\begin{quote}
After the accident a big crowd gathered there.
\end{quote}

A neuter noun.

\begin{cousinweb}
\textbf{\gu{ટોળું}} (\emph{ṭoḷũ}) comes from \emph{*ṭōla}, \textquotedblleft{}a division, a party\textquotedblright{}, a form CDIAL reconstructs; CDIAL lists Gujarati \emph{ṭoḷũ} there, glossed \textquotedblleft{}a flock of sheep\textquotedblright{} (CDIAL 5483). Today it is mostly a crowd of people that gathers at one spot, where \textbf{\gu{ભીડ}} is any crowd.
\end{cousinweb}

\begin{grammarlens}[title={What you've built}]
One more word for telling what happened.
\end{grammarlens}

\subsection*{Your turn}
\begin{itemize}
\raggedright
  \item \emph{Say it:} \emph{ṭoḷũ}
  \item \emph{Say it:} \emph{ṭoḷũ}, once more
  \item \emph{Say it:} \emph{akasmāt pachī tyā̃ lokonũ moṭũ ṭoḷũ bhegũ thayũ.} — After the accident a big crowd gathered there.
  \item \emph{Recall:} say the Gujarati for careful, then the Gujarati for upside down, then say \emph{ṭoḷũ} again
\end{itemize}

\begin{tcolorbox}[breakable,colback=teal!4,colframe=teal!35!black,title={Before you move on}]
What does \gu{ટોળું} mean? (\textquotedblleft{}A gathered crowd; a mob\textquotedblright{}.) Say it once more.
\end{tcolorbox}
\section[\texorpdfstring{\gu{ધક્કો} (dhakko)}{ધક્કો (dhakko)}]{\gu{ધક્કો} (dhakko) — a push}

\label{lesson:GU-C297-dhakko}



\begin{quote}
Before the new one: say the Gujarati for upside down, then the Gujarati for a gathered crowd.
\end{quote}

\subsection*{You'll want to know: \gu{ધક્કો}}
\textbf{\gu{ધક્કો}} — \emph{dhakko} — \textquotedblleft{}a push\textquotedblright{}.

\begin{quote}
\gu{ટોળામાં} \gu{કોઈએ} \gu{મને} \gu{ધક્કો} \gu{માર્યો}.
\end{quote}

\begin{quote}
\emph{ṭoḷāmā̃ koīe mane dhakko māryo.}
\end{quote}

\begin{quote}
Someone in the crowd pushed me.
\end{quote}

A masculine noun.

\begin{cousinweb}
\textbf{\gu{ધક્કો}} (\emph{dhakko}) comes from \emph{*dhakk}, \textquotedblleft{}push, strike\textquotedblright{}, a form CDIAL reconstructs, and CDIAL lists the Gujarati noun beside \textbf{\gu{ધકેલવું}}, to push (CDIAL 6701). A push is given with \textbf{\gu{મારવું}}: \textbf{\gu{ધક્કો} \gu{માર્યો}}.
\end{cousinweb}

\begin{grammarlens}[title={What you've built}]
One more word for telling what happened.
\end{grammarlens}

\subsection*{Your turn}
\begin{itemize}
\raggedright
  \item \emph{Say it:} \emph{dhakko}
  \item \emph{Say it:} \emph{dhakko}, once more
  \item \emph{Say it:} \emph{ṭoḷāmā̃ koīe mane dhakko māryo.} — Someone in the crowd pushed me.
  \item \emph{Recall:} say the Gujarati for upside down, then the Gujarati for a gathered crowd, then say \emph{dhakko} again
\end{itemize}

\begin{tcolorbox}[breakable,colback=teal!4,colframe=teal!35!black,title={Before you move on}]
What does \gu{ધક્કો} mean? (\textquotedblleft{}A push\textquotedblright{}.) Say it once more.
\end{tcolorbox}
